% ============ 第三章 III: 局部类域论: 上同调 ============
% 源: J.S. Milne, Class Field Theory (v4.03), Chapter III, pp. 72--106
% 翻译自 OCR 草稿 J.S. Milne - Class Field Theory.pdf_by_PaddleOCR-VL-1.6.md (行 3787--4584)
% 编号约定见 main.tex: 定理类环境按"节.序号"编号, 全部共享 theorem 计数器.
%   本章各节的编号条目均为定理类环境或叙述性条目, 序号连续, 无需 \setcounter 干预;
%   唯一例外是 §3 中的 3.3, 它是不带环境名的编号条目, 故在其后 \setcounter{theorem}{3}.
%   源书中的公式编号 (26)--(34) 用 \tag 如实保留, 使正文中的交叉引用得以落实.

\chapter{局部类域论: 上同调}\label{ch:III}\index{jubuleiyulun@局部类域论(local class field theory)}\index{shangtongdiao@上同调(cohomology)}

在本章中, 我们展开局部类域论的上同调进路. 我们首先证明局部 Artin 映射的存在性 (定理\ref{thm:iii.3.4}). 结合第 \ref{ch:I} 章 \S\S 1--3 的定理, 这就足以完成局部类域论主定理的证明. 对于那些跳过了 Lubin-Tate 理论的读者, 我们随后给出一个不依赖第 \ref{ch:I} 章 (\S 1 除外) 的主定理证明.

贯穿本章, ``局部域'' 指 ``非阿基米德局部域''. 一如既往, $K^{\mathrm{al}}$ 表示 $K$ 的一个代数闭包 (在 $K$ 的特征为 $p\neq 0$ 时指可分代数闭包), 而 ``$K$ 的扩张'' 指 ``$K^{\mathrm{al}}$ 中包含 $K$ 的子域''. 所有上同调群都用连续上链来计算 (见 II, \S 4). 对域的 Galois 扩张 $L/K$ (可能无限), 记
\[
 H^2(L/K)=H^2(\Gal(L/K),L^\times).
\]
眼下 $H^2(L/K)$ 只是一个上同调群, 但在下一章我们将会看到它作为 $L/K$ 的相对 Brauer 群有明确的解释.

% ==================================================================
\section{非分歧扩张的上同调}\label{sec:III.1}\index{feifenqikuozhang@非分歧扩张(unramified extension)}

\subsubsection*{单位群的上同调}

设 $K$ 为局部域.

\begin{proposition}\label{prop:iii.1.1}
设 $L/K$ 为 Galois 群为 $G$ 的有限非分歧扩张, 并设 $U_L$ 为 $L$ 中的单位群. 则
\[
 H_T^r(G,U_L)=0,\quad\text{所有 }r.
\]
\end{proposition}

\begin{proof}
若 $\pi$ 为 $L$ 的一个素元, 那么 $L^\times$ 的每个元素都可唯一写成 $\alpha=u\pi^m$ 的形式, 其中 $u\in U_L$, $m\in\Z$. 于是
\begin{equation}
 L^\times=U_L\cdot\pi^{\Z}\simeq U_L\times\Z.\tag{26}
\end{equation}
因为 $L$ 在 $K$ 上非分歧, 我们可以取 $\pi\in K$. 于是对 $\tau\in\Gal(L/K)$ 有 $\tau\alpha=\tau(u\pi^m)=(\tau u)\pi^m$, 因此当我们让 $G$ 平凡地作用在 $\pi^{\Z}\simeq\Z$ 上时, (26) 成为 $G$-模的一个分解. 所以 $H^r(G,U_L)$ 是 $H^r(G,L^\times)$ 的一个直和因子 (见 II, \ref{prop:ii.1.25}). 由于依 Hilbert 定理 90 (II, \ref{prop:ii.1.22}) 有 $H^1(G,L^\times)=0$, 这表明 $H^1(G,U_L)=0$. 因为 $G$ 是循环群, 为完成定理的证明, 只需 (依 II, \ref{prop:ii.3.4}) 证明 $H_T^0(G,U_L)=0$. 这一点由下一个命题完成.
\end{proof}

\begin{proposition}\label{prop:iii.1.2}
设 $L/K$ 为有限非分歧扩张. 则范映射 $\Nm_{L/K}\colon U_L\to U_K$ 是满射.
\end{proposition}

我们先需要一些引理. 设 $l$ 与 $k$ 为 $L$ 与 $K$ 的剩余类域. 因为 $L/K$ 非分歧, $G$ 在 $\mathcal{O}_L$ 上的作用定义了一个同构 $G\simeq\Gal(l/k)$.

\begin{lemma}\label{lem:iii.1.3}
对 $m>0$, 令 $U_L^{(m)}=1+\mathfrak{m}_L^m$. 则作为 $G$-模有
\[
 U_L/U_L^{(1)}\xrightarrow{\simeq} l^\times
\]
\[
 U_L^{(m)}/U_L^{(m+1)}\xrightarrow{\simeq} l .
\]
\end{lemma}

\begin{proof}
设 $\pi$ 为 $K$ 的一个素元. 它在 $L$ 中仍是素元, 且
\[
 U_L^{(m)}=\{1+a\pi^m\mid a\in\mathcal{O}_L\}.
\]
映射
\[
 u\mapsto u\mod\mathfrak{m}_L\colon U_L\to l^\times
\]
\[
 1+a\pi^m\mapsto a\mod\mathfrak{m}_L\colon U_L^{(m)}\to l
\]
诱导出所需的同构.
\end{proof}

\begin{lemma}\label{lem:iii.1.4}
对所有 $r$ 有 $H_T^r(G,l^\times)=0$. 特别地, 范映射 $l^\times\to k^\times$ 是满射.
\end{lemma}

\begin{proof}
依 Hilbert 定理 90 (II, \ref{prop:ii.1.22}), $H^1(G,l^\times)=0$, 又因为 $l^\times$ 有限, 其 Herbrand 商 $h(l^\times)=1$ (见 II, \ref{prop:ii.3.8}). 因此 $H^2(G,l^\times)=0$, 依 (II, \ref{prop:ii.3.4}) 这推出所有的群都是零.
\end{proof}

\begin{lemma}\label{lem:iii.1.5}
群 $H_T^r(G,l)=0$ 对所有 $r$ 成立. 特别地, 迹映射 $l\to k$ 是满射.
\end{lemma}

\begin{proof}
在 (II, \ref{prop:ii.1.24}) 中对 $r>0$ 证明了这一点, 依 (II, \ref{prop:ii.3.4}) 这蕴含它对所有 $r$ 成立.
\end{proof}

\begin{proof}[命题\ref{prop:iii.1.2} 的证明]
有交换图表 ($m>0$)
\[
\begin{tikzcd}
	{U_L} & {l^\times} \\
	{U_K} & {k^\times}
	\arrow[from=1-1, to=1-2]
	\arrow["\Nm"', from=1-1, to=2-1]
	\arrow["\Nm", from=1-2, to=2-2]
	\arrow[from=2-1, to=2-2]
\end{tikzcd}
\qquad
\begin{tikzcd}
	{U_L^{(m)}} & l \\
	{U_K^{(m)}} & k
	\arrow[from=1-1, to=1-2]
	\arrow["\Nm"', from=1-1, to=2-1]
	\arrow["\Tr", from=1-2, to=2-2]
	\arrow[from=2-1, to=2-2]
\end{tikzcd}
\]
取 $u\in U_K$. 因为范映射 $l^\times\to k^\times$ 是满射, 存在 $v_0\in U_L$ 使得 $\Nm(v_0)$ 与 $u$ 在 $k^\times$ 中有相同的像, 即 $u/\Nm(v_0)\in U_K^{(1)}$. 因为迹映射 $l\to k$ 是满射, 存在 $v_1\in U_L^{(1)}$ 使得 $\Nm(v_1)\equiv u/\Nm(v_0)\mod U_K^{(2)}$. 如此继续下去, 我们得到一列元素 $v_0,v_1,v_2,v_3,\ldots$, $v_i\in U_L^{(i)}$, 使得 $u/\Nm(v_0\cdots v_i)\in U_K^{(i+1)}$. 令 $v=\lim_{m\to\infty}\prod_{j=0}^{m}v_j$. 则 $u/\Nm(v)\in\bigcap U_K^{(i)}=\{1\}$.
\end{proof}

\begin{corollary}\label{cor:iii.1.6}
设 $L/K$ 为 Galois 群为 $G$ 的无限非分歧扩张. 则对 $r>0$ 有 $H^r(G,U_L)=0$ (连续上链).
\end{corollary}

\begin{proof}
域 $L$ 是 $K$ 的有限扩张 $K'$ 的并, 于是 (见 II, \ref{prop:ii.4.2})
\[
 H^r(\Gal(L/K),U_L)\simeq\varinjlim_{K'}H^r(\Gal(K'/K),U_{K'})=0.
\]
\end{proof}

\subsubsection*{不变映射}

设 $L$ 为 $K$ 的一个非分歧扩张 (可能无限), 并设 $G=\Gal(L/K)$.

由于 $H^2(G,U_L)=0=H^3(G,U_L)$, 短正合列
\[
 0\to U_L\to L^\times\xrightarrow{\ord_L}\Z\to 0
\]
的上同调序列给出一个同构
\[
 H^2(G,L^\times)\xrightarrow[\simeq]{\quad H^2(\ord_L)\quad}H^2(G,\Z).
\]
群 $H^r(G,\Q)$ 对 $r>0$ 是挠群 (见 II, \ref{cor:ii.4.3}) 并且是唯一可除的 (因为 $\Q$ 如此), 故为零. 因此短正合列
\[
 0\to\Z\to\Q\to\Q/\Z\to 0
\]
(带平凡的 $G$-作用) 的上同调序列给出一个同构
\[
 H^1(G,\Q/\Z)\xrightarrow{\delta}H^2(G,\Z).
\]
回顾
\[
 H^1(G,\Q/\Z)\simeq\Hom_{\mathrm{cts}}(G,\Q/\Z),
\]
并且 $G$ 有一个典范的拓扑生成元, 即 Frobenius 元素 $\sigma=\Frob_{L/K}$. 复合映射
\[
 H^2(L/K)\xrightarrow[\simeq]{\ord_L}H^2(G,\Z)\xleftarrow[\simeq]{\delta}H^1(G,\Q/\Z)\simeq\Hom_{\mathrm{cts}}(G,\Q/\Z)\xrightarrow{f\mapsto f(\sigma)}\Q/\Z
\]
称为\emph{不变映射}
\[
 \inv_{L/K}\colon H^2(L/K)\to\Q/\Z.
\]

考虑一个域扩张塔
\[
 E\supset L\supset K,
\]
其中 $E$ 与 $L$ 都在 $K$ 上非分歧 (从而 Galois). 那么
\[
\begin{tikzcd}
	{H^2(L/K)} & {\Q/\Z} \\
	{H^2(E/K)} & {\Q/\Z}
	\arrow["{\inv_{L/K}}", from=1-1, to=1-2]
	\arrow["\Inf"', from=1-1, to=2-1]
	\arrow[equal, from=1-2, to=2-2]
	\arrow["{\inv_{E/K}}"', from=2-1, to=2-2]
\end{tikzcd}
\]
交换, 因为 $\inv$ 的定义中所有的映射都与 $\Inf$ 相容.

\begin{theorem}\label{thm:iii.1.7}
存在唯一一个同构
\[
 \inv_K\colon H^2(K^{\mathrm{un}}/K)\to\Q/\Z,
\]
具有以下性质: 对每个在 $K$ 上 $n$ 次有限的 $L\subset K^{\mathrm{un}}$, $\inv_K$ 诱导出同构
\[
 \inv_{L/K}\colon H^2(L/K)\to\tfrac{1}{[L:K]}\Z/\Z.
\]
\end{theorem}

\begin{proof}
这是上面讨论的直接推论.
\end{proof}

\begin{proposition}\label{prop:iii.1.8}
设 $L$ 为 $K$ 的 $n$ 次有限扩张, 并设 $K^{\mathrm{un}}$ 与 $L^{\mathrm{un}}$ 为 $K$ 与 $L$ 的最大非分歧扩张. 则下列图表交换:
\begin{equation}
\begin{tikzcd}[column sep=small]
	{H^2(K^{\mathrm{un}}/K)} & {H^2(L^{\mathrm{un}}/L)} \\
	{\Q/\Z} & {\Q/\Z}
	\arrow["\Res", from=1-1, to=1-2]
	\arrow["{\inv_K}"', from=1-1, to=2-1]
	\arrow["{\inv_L}", from=1-2, to=2-2]
	\arrow["n"', from=2-1, to=2-2]
\end{tikzcd}\tag{27}
\end{equation}
\end{proposition}

\begin{proof}
局部域的最大非分歧扩张可由添加所有 $m$ 次单位根得到, 其中 $m$ 不被剩余特征整除. 因此 $L^{\mathrm{un}}=L\cdot K^{\mathrm{un}}$, 从而映射
\[
 \tau\mapsto\tau|K^{\mathrm{un}}\colon \Gal(L^{\mathrm{un}}/L)\to\Gal(K^{\mathrm{un}}/K)
\]
是单射. (27) 中记作 $\Res$ 的映射由相容的同态
\[
\begin{tikzcd}
	{\Gal(K^{\mathrm{un}}/K)} & {\Gal(L^{\mathrm{un}}/L)} \\
	{K^{\mathrm{un}\times}} & {L^{\mathrm{un}\times}}
	\arrow[from=2-1, to=2-2]
\end{tikzcd}
\]
所定义 (II, \ref{ex:ii.1.27}b); 其中第一行由 $\tau\mapsto\tau|K^{\mathrm{un}}$ 的反向给出.

令 $\Gamma_K=\Gal(K^{\mathrm{un}}/K)$ 与 $\Gamma_L=\Gal(L^{\mathrm{un}}/L)$, 并考虑图表
\[
\begin{tikzcd}[column sep=small, row sep=large]
	{H^2(K^{\mathrm{un}}/K)} & {H^2(\Gamma_K,\Z)} & {H^1(\Gamma_K,\Q/\Z)} & {\Q/\Z} \\
	{H^2(L^{\mathrm{un}}/L)} & {H^2(\Gamma_L,\Z)} & {H^1(\Gamma_L,\Q/\Z)} & {\Q/\Z}
	\arrow["{\ord_K}"{description}, "\simeq"{above}, from=1-1, to=1-2]
	\arrow["\delta"{description}, "\simeq"{above}, from=1-3, to=1-2]
	\arrow["{g\mapsto g(\sigma_K)}", from=1-3, to=1-4]
	\arrow["{\ord_L}"{description}, "\simeq"{above}, from=2-1, to=2-2]
	\arrow["\delta"{description}, "\simeq"{above}, from=2-3, to=2-2]
	\arrow["{g\mapsto g(\sigma_L)}"', from=2-3, to=2-4]
	\arrow["\Res"', from=1-1, to=2-1]
	\arrow["{e\,\Res}", from=1-2, to=2-2]
	\arrow["\Res"', from=1-3, to=2-3]
	\arrow["f", from=1-4, to=2-4]
\end{tikzcd}
\]
这里 $e$ 是 $L/K$ 的分歧指数, $f$ 是剩余类次数. 第一个方块来自交换方块
\[
\begin{tikzcd}
	{K^{\mathrm{un}\times}} & \Z \\
	{L^{\mathrm{un}\times}} & \Z
	\arrow["{\ord_K}", from=1-1, to=1-2]
	\arrow[from=1-1, to=2-1]
	\arrow["e", from=1-2, to=2-2]
	\arrow["{\ord_L}"', from=2-1, to=2-2]
\end{tikzcd}
\]
第二个方块表达的事实是限制映射与边缘映射可交换. 除因子 ``$e$'' 外, 第三个方块是
\[
\begin{tikzcd}
	{\Hom(\Gamma_K,\Q/\Z)} & {\Q/\Z} \\
	{\Hom(\Gamma_L,\Q/\Z)} & {\Q/\Z}
	\arrow["{g\mapsto g(\sigma_K)}", from=1-1, to=1-2]
	\arrow["{g\mapsto g|\Gamma_L}"', from=1-1, to=2-1]
	\arrow["f", from=1-2, to=2-2]
	\arrow["{g\mapsto g(\sigma_L)}"', from=2-1, to=2-2]
\end{tikzcd}
\]
Frobenius 元素 $\sigma_K$ 与 $\sigma_L$ 由它们在剩余类域上分别诱导 $x\mapsto x^q$ 与 $x\mapsto x^{q^f}$ 这一事实所确定, 其中 $q=|k|$, $q^f=|l|$, 因此 $\sigma_L|K^{\mathrm{un}}=\sigma_K^f$. 现在该方块交换是清楚的, 又因为 $n=ef$, 这就证明了命题.
\end{proof}

\subsubsection*{局部 Artin 映射}

设 $L$ 为 $K$ 的 Galois 群为 $G$ 的有限非分歧扩张, 并令 $n=[L:K]$. \emph{局部基本类} $u_{L/K}$ 是 $H^2(L/K)$ 中被不变映射
\[
 \inv_{L/K}\colon H^2(L/K)\xrightarrow{\simeq}\tfrac{1}{[L:K]}\Z/\Z
\]
映到生成元 $\tfrac{1}{[L:K]}$ 的那个元素. 偶对 $(G,L^\times)$ 满足 Tate 定理 (II, \ref{thm:ii.3.11}) 的假设,\footnote{对 $H^1$ 的条件, 用 Hilbert 定理 90; 对 $H^2$ 的条件, 用不变映射.} 于是与基本类 $u_{L/K}$ 作杯积定义了同构
\[
 H_T^r(G,\Z)\to H_T^{r+2}(G,L^\times)
\]
(对所有 $r\in\Z$). 当 $r=-2$ 时, 这成为
\[
\begin{tikzcd}
	{H^{-2}(G,\Z)} & {H^0(G,L^\times)} \\
	G & {K^\times/\Nm(L^\times)}
	\arrow["\simeq", from=1-1, to=1-2]
	\arrow[equal, from=1-1, to=2-1]
	\arrow[equal, from=1-2, to=2-2]
\end{tikzcd}
\]
我们现在明确地计算这个映射.

$K$ 的一个素元 $\pi$ 也是 $L$ 的素元, 并给出一个 $G$-模的分解
\[
 L^\times=U_L\cdot\pi^{\Z}\simeq U_L\times\Z.
\]
于是
\[
 H^r(G,L^\times)\simeq H^r(G,U_L)\oplus H^r(G,\pi^{\Z}).
\]
取 $G$ 的一个生成元 $\sigma$ (例如 Frobenius 生成元), 并令
\[
 f\in H^1(G,\Q/\Z)\simeq\Hom(G,\Q/\Z)
\]
为满足对所有 $i$ 有 $f(\sigma^i)=\tfrac{i}{n}\mod\Z$ 的元素. 它生成 $H^1(G,\Q/\Z)$.

从正合列
\[
 0\to\Z\to\Q\to\Q/\Z\to 0
\]
以及 $H^r(G,\Q)=0$ (对所有 $r$) 这一事实, 我们得到一个同构
\[
 \delta\colon H^1(G,\Q/\Z)\to H^2(G,\Z).
\]
按照 (II, \ref{rem:ii.1.21}) 中对 $\delta$ 的描述, 为构造 $\delta f$, 我们首先把 $f$ 提升为一个 $1$-上链 $\tilde f\colon G\to\Q$. 我们取 $\tilde f$ 为映射 $\sigma^i\mapsto\tfrac{i}{n}$, 其中 $0\le i<n$. 那么
\[
 d\tilde f(\sigma^i,\sigma^j)=\sigma^i\tilde f(\sigma^j)-\tilde f(\sigma^{i+j})+\tilde f(\sigma^i)=
 \begin{cases}
 0 & \text{若 } i+j\le n-1,\\
 1 & \text{若 } i+j>n-1.
 \end{cases}
\]
当把 $\Z$ 与 $L^\times$ 的子群 $\pi^{\Z}$ 等同起来时, 我们发现基本类 $u_{L/K}\in H^2(G,L^\times)$ 由上循环
\[
 \varphi(\sigma^i,\sigma^j)=
 \begin{cases}
 1 & \text{若 } i+j\le n-1,\\
 \pi & \text{若 } i+j>n-1
 \end{cases}
\]
表示.

从正合列
\[
 0\to I\to\Z[G]\to\Z\to 0
\]
\[
 0\to L^\times\to L^\times(\varphi)\to I\to 0
\]
(见 II, \ref{thm:ii.3.11} 的证明) 我们得到边缘映射
\[
 H^{-2}(G,\Z)\to H^{-1}(G,I)
\]
\[
 H^{-1}(G,I)\to H^0(G,L^\times),
\]
因为 $\Z[G]$ 与 $L^\times(\varphi)$ 有平凡的上同调, 它们是同构. 这里 $L^\times(\varphi)$ 是 $\varphi$ 的分裂模 $L^\times\oplus\bigoplus_{\sigma\in G,\sigma\neq 1}\Z x_\sigma$.

最后, $H^{-2}(G,\Z)\overset{\mathrm{def}}{=}H_1(G,\Z)\simeq G$ (见 II, \ref{prop:ii.2.7}).

\begin{proposition}\label{prop:iii.1.9}
在复合映射
\[
\begin{tikzcd}
	{H^{-2}(G,\Z)} & {H^0(G,L^\times)} \\
	G & {K^\times/\Nm(L^\times)}
	\arrow["\simeq", from=1-1, to=1-2]
	\arrow[equal, from=1-1, to=2-1]
	\arrow[equal, from=1-2, to=2-2]
\end{tikzcd}
\]
下, Frobenius 元素 $\sigma\in G$ 映到 $\pi$ 在 $K^\times/\Nm(L^\times)$ 中的类.
\end{proposition}

注意, 因为 $K$ 中所有单位都是来自 $L^\times$ 的范 (见 \ref{prop:iii.1.2}), $\pi\mod\Nm(L^\times)$ 的类与素元 $\pi$ 的选取无关. 另一方面, $G$-模 $L^\times(\varphi)$ 以及该映射依赖于 $G$ 的生成元 $\sigma$ 的选取.

\begin{proof}
从同构 $H^{-2}(G,\Z)\simeq G$ 的构造我们看到, $\sigma$ 在边缘映射 $H^{-2}(G,\Z)\to H^{-1}(G,I_G)\subset I_G/I_G^2$ 下的像由 $\sigma-1$ 表示.

边缘映射 $H^{-1}(G,I_G)\to H^0(G,L^\times)$ 是由下列图表经蛇形引理给出的那个映射 (我们把 $I_G$ 记作 $I$):
\[
\begin{tikzcd}[column sep=small]
	0 & {L^\times} & {L^\times(\varphi)} & I & 0 \\
	0 & {L^\times} & {L^\times(\varphi)} & I & 0
	\arrow[from=1-1, to=1-2]
	\arrow[from=1-2, to=1-3]
	\arrow[from=1-3, to=1-4]
	\arrow[from=1-4, to=1-5]
	\arrow["{\Nm_G}"', from=1-2, to=2-2]
	\arrow["{\Nm_G}"', from=1-3, to=2-3]
	\arrow["{\Nm_G}"', from=1-4, to=2-4]
	\arrow[from=2-1, to=2-2]
	\arrow[from=2-2, to=2-3]
	\arrow[from=2-3, to=2-4]
	\arrow[from=2-4, to=2-5]
\end{tikzcd}
\]
连接两行的竖直映射是 $\Nm_G=\sum_{i=0}^{n-1}\sigma^i$. 元素 $(\sigma-1)+I^2$ 是 $x_\sigma+I\cdot L^\times(\varphi)$ 在 $L^\times(\varphi)_G$ 中的像, 而 $\Nm_G(x_\sigma+I\cdot L^\times(\varphi))$ 是下列元素之和:
\[
\begin{aligned}
 x_\sigma &= x_\sigma\\
 \sigma x_\sigma &= x_{\sigma^2}-x_\sigma+\varphi(\sigma,\sigma)\\
 \sigma^2 x_\sigma &= x_{\sigma^3}-x_{\sigma^2}+\varphi(\sigma,\sigma^2)\\
 &\;\;\vdots\\
 \sigma^{n-1}x_\sigma &= x_1-x_{\sigma^{n-1}}+\varphi(\sigma,\sigma^{n-1}).
\end{aligned}
\]
把这些加起来, 记住在 $L^\times(\varphi)$ 的因子 $L^\times$ 上 ``$+$'' 实际上是 ``$\cdot$'', 我们得到
\[
 \Nm_G(x_\sigma)=\prod_{i=1}^{n-1}\varphi(\sigma,\sigma^i)=\pi.
\]
证明完成.
\end{proof}

\begin{remark}\label{rem:iii.1.10}
上面这个命题\ref{prop:iii.1.9} 的证明使用了 Tate 关于同构 $H^r(G,\Z)\to H^{r+2}(G,C)$ 的原始定义, 它比其他文献中用杯积描述该映射的证明更简单.
\end{remark}

% ==================================================================
\section{分歧扩张的上同调}\label{sec:III.2}\index{fenqikuozhang@分歧扩张(ramified extension)}

由于 Hilbert 定理 90 (II, \ref{prop:ii.1.22}), 对任意 Galois 扩张塔 $E\supset L\supset K$ 都有一个正合列
\[
 0\to H^2(L/K)\xrightarrow{\Inf}H^2(E/K)\xrightarrow{\Res}H^2(E/L)
\]
(见 II, \ref{ex:ii.1.36}).

下一个定理把定理\ref{thm:iii.1.7} 推广到分歧扩张.

\begin{theorem}\label{thm:iii.2.1}
对每个局部域 $K$, 存在一个典范同构
\[
 \inv_K\colon H^2(K^{\mathrm{al}}/K)\to\Q/\Z.
\]
设 $L$ 为 $K$ 的 $n<\infty$ 次 Galois 扩张. 那么图表
\begin{equation}
\begin{tikzcd}[column sep=small]
	0 & {H^2(L/K)} & {H^2(K^{\mathrm{al}}/K)} & {H^2(K^{\mathrm{al}}/L)} \\
	0 & {\frac{1}{n}\Z/\Z} & {\Q/\Z} & {\Q/\Z}
	\arrow[from=1-1, to=1-2]
	\arrow[from=1-2, to=1-3]
	\arrow["\Res", from=1-3, to=1-4]
	\arrow["{\inv_{L/K}}"', from=1-2, to=2-2]
	\arrow["{\inv_K}"', from=1-3, to=2-3]
	\arrow["{\inv_L}", from=1-4, to=2-4]
	\arrow[from=2-1, to=2-2]
	\arrow[from=2-2, to=2-3]
	\arrow["n"', from=2-3, to=2-4]
\end{tikzcd}\tag{28}
\end{equation}
交换, 从而定义了一个同构
\[
 \inv_{L/K}\colon H^2(L/K)\to\tfrac{1}{n}\Z/\Z.
\]
\end{theorem}

证明将占据本节余下的大部分篇幅.

\begin{lemma}\label{lem:iii.2.2}
若 $L/K$ 是 $n$ 次 Galois 扩张, 则 $H^2(L/K)$ 包含一个典范同构于 $\tfrac{1}{n}\Z/\Z$ 的子群.
\end{lemma}

\begin{proof}
考虑图表
\[
\begin{tikzcd}[column sep=small]
	0 & {\Ker(\Res)} & {H^2(K^{\mathrm{un}}/K)} & {H^2(L^{\mathrm{un}}/L)} \\
	0 & {H^2(L/K)} & {H^2(K^{\mathrm{al}}/K)} & {H^2(K^{\mathrm{al}}/L)}
	\arrow[from=1-1, to=1-2]
	\arrow[from=1-2, to=1-3]
	\arrow["\Res", from=1-3, to=1-4]
	\arrow[from=1-2, to=2-2]
	\arrow["\Inf"', from=1-3, to=2-3]
	\arrow["\Inf", from=1-4, to=2-4]
	\arrow[from=2-1, to=2-2]
	\arrow[from=2-2, to=2-3]
	\arrow["\Res"', from=2-3, to=2-4]
\end{tikzcd}
\]
由于两个膨胀映射都是单射, 第一个竖直映射也是单射, 而 \ref{prop:iii.1.8} 表明顶行限制映射的核典范同构于 $\tfrac{1}{n}\Z/\Z$.
\end{proof}

为完成定理\ref{thm:iii.2.1} 的证明, 只需证明映射 $\tfrac{1}{n}\Z/\Z\hookrightarrow H^2(L/K)$ 是同构. 有两种不同的途径来证明这一点. 在下一章关于 Brauer 群的内容中, 我们将证明 $H^2(K^{\mathrm{un}}/K)\simeq H^2(K^{\mathrm{al}}/K)$ ($K$ 上每个中心单代数都被一个非分歧扩张分裂; 见 IV, \S 4), 于是上面图表中的竖直映射都是同构. 第二种证明, 也就是我们现在给出的证明, 则表明 $|H^2(L/K)|\le n$.

\begin{lemma}\label{lem:iii.2.3}
设 $L$ 为 $K$ 的 Galois 群为 $G$ 的有限 Galois 扩张. 则存在 $\mathcal{O}_L$ 的一个在 $G$ 下稳定的开子群 $V$, 使得对所有 $r>0$ 有 $H^r(G,V)=0$.
\end{lemma}

\begin{proof}
设 $\{x_\tau\mid\tau\in G\}$ 为 $L$ 在 $K$ 上的一个正规基 (见 FT, 5.18). 诸 $x_\tau$ 在 $\mathcal{O}_K$ 中有一个公分母 $d$ (见 ANT, 2.6). 把每个 $x_\tau$ 换成 $d\cdot x_\tau$ 之后, 可设它们都落在 $\mathcal{O}_L$ 中. 取 $V=\sum\mathcal{O}_K x_\tau$. 因为正规基在 $G$ 下稳定, $V$ 也在 $G$ 下稳定. 设 $\pi$ 为 $\mathcal{O}_L$ 的一个素元. 那么对某个 $m>0$ 有 $V\supset\pi^m\mathcal{O}_L$, 这表明 $V$ 是开的 (它是 $\pi^m\mathcal{O}_L$ 的陪集之并). 最后, 作为 $G$-模
\[
 V\simeq\mathcal{O}_K[G]\simeq\Ind^G\mathcal{O}_K,
\]
于是对所有 $r>0$ 有 $H^r(G,V)=0$ (II, \ref{cor:ii.1.12}).
\end{proof}

\begin{lemma}\label{lem:iii.2.4}
设 $L$, $K$, $G$ 如上个引理. 则存在 $U_L$ 的一个在 $G$ 下稳定的开子群 $V$, 使得对所有 $r>0$ 有 $H^r(G,V)=0$.
\end{lemma}

\begin{proof}
我只对特征零的 $K$ 给出证明.\footnote{特征 $p$ 情形的证明见 \cite{Serre1967b}, p.~134.} 幂级数
\[
 e^x=1+x+\cdots+x^n/n!+\cdots
\]
在 $\ord(x)>\ord(p)/(p-1)$ 时收敛 (ANT, 7.29). 它定义了 $L$ 中 $0$ 的一个开邻域到 $L^\times$ 中 $1$ 的一个开邻域上的同构, 其逆映射为
\[
 \log(x)=(x-1)-(x-1)^2/2+(x-1)^3/3-\cdots.
\]
两个映射都与 $G$ 的作用可交换 (因为 $G$ 作用是连续的). 若 $V'$ 是 (\ref{lem:iii.2.3}) 中那样的 $0$ 的开邻域, 则 $\pi^M V'$ 有同样的性质, 于是我们可以取 $V=\exp(\pi^M V')$, 其中 $M$ 取得足够大, 使得指数函数在 $\pi^M V'$ 上有定义并且是同构.
\end{proof}

\begin{lemma}\label{lem:iii.2.5}
设 $L/K$ 为 $n$ 次循环扩张; 则 $h(U_L)=1$ 且 $h(L^\times)=n$.
\end{lemma}

\begin{proof}
设 $V$ 为 $U_L$ 的一个开子群, 满足对所有 $r$ 有 $H^r(G,V)=0$. 因为 $U_L$ 是紧的, 商群 $U_L/V$ 有限, 于是依 (II, \ref{cor:ii.3.9}) 有 $h(U_L)=h(V)=1$. 现在 $h(L^\times)=h(U_L)\cdot h(\Z)=h(\Z)$, 且
\[
 h(\Z)=\bigl|H_T^0(G,\Z)\bigr|/\bigl|H^1(G,\Z)\bigr|=\bigl|H_T^0(G,\Z)\bigr|=\bigl|\Z/n\Z\bigr|=n.
\]
\end{proof}

\begin{lemma}\label{lem:iii.2.6}
设 $L$ 为 $K$ 的 $n$ 阶有限 Galois 扩张, 则 $H^2(L/K)$ 的阶为 $n$.
\end{lemma}

\begin{proof}
我们知道 $H^2(L/K)$ 的阶被 $n$ 整除, 并且当 $L/K$ 循环时它等于 $n$. 我们对 $[L:K]$ 作归纳来证明引理. 因为群 $\Gal(L/K)$ 是可解的 (ANT, 7.59), 存在 Galois 扩张 $K'/K$ 满足 $L\supsetneq K'\supsetneq K$. 从正合列
\[
 0\to H^2(K'/K)\to H^2(L/K)\to H^2(L/K')
\]
我们看出
\[
 \bigl|H^2(L/K)\bigr|\le\bigl|H^2(K'/K)\bigr|\cdot\bigl|H^2(L/K')\bigr|=n.
\]
\end{proof}

\begin{proof}[定理\ref{thm:iii.2.1} 的证明]
从 (\ref{lem:iii.2.2}) 证明中的图表我们看到, 对 $K$ 的每个有限 Galois 扩张 $L$, $H^2(K^{\mathrm{al}}/K)$ 的子群 $H^2(L/K)$ 包含在 $H^2(K^{\mathrm{un}}/K)$ 中. 由于 $H^2(K^{\mathrm{al}}/K)=\bigcup H^2(L/K)$, 这证明了膨胀映射 $H^2(K^{\mathrm{un}}/K)\to H^2(K^{\mathrm{al}}/K)$ 是同构. 把它求逆再与定理\ref{thm:iii.1.7} 的不变映射 $\inv_K\colon H^2(K^{\mathrm{un}}/K)\to\Q/\Z$ 复合即可.
\end{proof}

事实上, 定理\ref{thm:iii.2.1} 中的图表即使当 $L/K$ 不是 Galois 扩张时也交换:
\begin{equation}
 \inv_L\circ\Res=[L:K]\inv_K.\tag{29}
\end{equation}
(在这种情形下, 我们把 $H^2(L/K)$ 定义为限制映射的核.)

我断言
\begin{equation}
 \inv_K\circ\Cor=\inv_L.\tag{30}
\end{equation}
因为 $\Res$ 是满射 (见图表 (28)), 为证明这一点只需证明
\[
 \inv_K\circ\Cor\circ\Res=\inv_L\circ\Res.
\]
但依 (II, \ref{prop:ii.1.30}) 有 $\Cor\circ\Res=[L:K]$, 于是这就是前面的等式 (29).

\subsubsection*{基本类}

设 $L$ 为 $K$ 的 Galois 群为 $G$ 的有限 Galois 扩张. 如同非分歧的情形, 我们定义 $L/K$ 的\emph{基本类} $u_{L/K}$ 为 $H^2(L/K)$ 中满足
\[
 \inv_{L/K}(u_{L/K})=\frac{1}{[L:K]}\bmod\Z
\]
的元素, 等价地, 满足
\[
 \inv_K(u_{L/K})=\frac{1}{[L:K]}\bmod\Z
\]
的元素.

\begin{lemma}\label{lem:iii.2.7}
设 $L\supset E\supset K$, 其中 $L/K$ 有限 Galois. 则
\begin{equation}
 \Res(u_{L/K})=u_{L/E},\tag{31}
\end{equation}
\begin{equation}
 \Cor(u_{L/E})=[E:K]u_{L/K}.\tag{32}
\end{equation}
此外, 若 $E/K$ 也是 Galois 扩张, 则
\begin{equation}
 \Inf(u_{E/K})=[L:E]u_{L/K}.\tag{33}
\end{equation}
\end{lemma}

\begin{proof}
考虑
\[
\begin{tikzcd}[column sep=small]
	{H^2(K^{\mathrm{al}}/K)} & {H^2(K^{\mathrm{al}}/E)} & {H^2(K^{\mathrm{al}}/L)} \\
	{\Q/\Z} & {\Q/\Z} & {\Q/\Z}
	\arrow["\Res", from=1-1, to=1-2]
	\arrow["\Res", from=1-2, to=1-3]
	\arrow["{\inv_K}"', "\simeq"', from=1-1, to=2-1]
	\arrow["{\inv_E}"', "\simeq"', from=1-2, to=2-2]
	\arrow["{\inv_L}", "\simeq"', from=1-3, to=2-3]
	\arrow["{[E:K]}"', from=2-1, to=2-2]
	\arrow["{[L:E]}"', from=2-2, to=2-3]
\end{tikzcd}
\]
把核--余核引理 (II, \ref{lem:ii.A.2}) 用于两行 (并利用该引理的函子性), 我们得到一个交换图表
\[
\begin{tikzcd}[column sep=small]
	0 & {H^2(E/K)} & {H^2(L/K)} & {H^2(L/E)} \\
	0 & {\frac{1}{[E:K]}\Z/\Z} & {\frac{1}{[L:K]}\Z/\Z} & {\frac{1}{[L:E]}\Z/\Z}
	\arrow[from=1-1, to=1-2]
	\arrow["\Inf", from=1-2, to=1-3]
	\arrow["\Res", from=1-3, to=1-4]
	\arrow["{\inv_{E/K}}"', from=1-2, to=2-2]
	\arrow["{\inv_{L/K}}"', from=1-3, to=2-3]
	\arrow["{\inv_{L/E}}", from=1-4, to=2-4]
	\arrow[from=2-1, to=2-2]
	\arrow["\id"{description}, from=2-2, to=2-3]
	\arrow["{[E:K]}"', from=2-3, to=2-4]
\end{tikzcd}
\]
这里 $H^2(E/K)$ 定义为 $\Res\colon H^2(K^{\mathrm{al}}/K)\to H^2(K^{\mathrm{al}}/E)$ 的核. 等式 (31) 表达的是第二个方块交换. 当 $E/K$ 为 Galois 扩张时, 等式 (33) 表达的是第一个方块交换.

等式 (32) 可由比较
\[
\begin{aligned}
 \inv_K(\Cor_{E/K}(u_{L/E}))&\overset{(30)}{=}\inv_E(u_{L/E})=\frac{1}{[L:E]}\bmod\Z,\\
 \inv_K(u_{L/K})&=\frac{1}{[L:K]}\bmod\Z
\end{aligned}
\]
得到.
\end{proof}

\begin{remark}\label{rem:iii.2.8}
对每个局部域的有限扩张 $L/K$ 都附有群 $K^\times/\Nm(L^\times)$ 与 $H^2(L/K)$. 当 $L/K$ 循环时它们同构, 否则不同构. 第一个群总是同构于 $\Gal(M/K)$, 其中 $M$ 是 $L/K$ 中最大的 Abel 子扩张 (见 \ref{thm:iii.3.4}, \ref{thm:iii.3.5}), 而第二个群总是 $[L:K]$ 阶循环群. 于是, 当 $L/K$ Abel 但不循环时, 这两个群有相同的阶但不同构; 而当 $L/K$ 不 Abel 时, 它们的阶不同.
\end{remark}

% ==================================================================
\section{局部 Artin 映射}\label{sec:III.3}\index{jubuartinying@局部 Artin 映射(local Artin map)}\index{jubuhuxiang@局部互反映射(local reciprocity map)}

偶对 $(G,L^\times)$ 满足 Tate 定理 (II, \ref{thm:ii.3.11}) 的假设, 于是我们证明了下述结果.

\begin{theorem}\label{thm:iii.3.1}
对局部域的每个有限 Galois 扩张 $L/K$ 以及 $r\in\Z$, 由 $x\mapsto x\cup u_{L/K}$ 定义的同态
\[
 H_T^r(\Gal(L/K),\Z)\to H_T^{r+2}(\Gal(L/K),L^\times)
\]
是同构. 当 $r=-2$ 时, 这成为同构
\[
 G^{\mathrm{ab}}\simeq K^\times/\Nm_{L/K}(L^\times).
\]

若 $L\supset E\supset K$ 且 $L/K$ Galois, 则定理中的同态与 $\Res$ 及 $\Cor$ 都可交换. 例如, 为证明关于 $\Res$ 的陈述, 我们必须证明
\[
 \Res(x\cup u_{L/K})=\Res(x)\cup u_{L/E}
\]
对所有 $x\in H_T^r(\Gal(L/K),L^\times)$ 成立. 但杯积的一个标准性质 (II, \ref{prop:ii.1.39}c) 是
\[
 \Res(x\cup u_{L/K})=\Res(x)\cup\Res(u_{L/K}),
\]
于是这由 (31) 推出. 关于 $\Cor$ 的证明类似.

若 $L\supset E\supset K$, 其中 $L/K$ 与 $E/K$ 都 Galois, 且 $x\in H^r(\Gal(E/K),E^\times)$, $r\ge 1$, 则
\[
 \Inf(x\cup u_{E/K})=[L:E]\Inf(x)\cup u_{L/K}.
\]
同样, 这由杯积的一个标准性质 (\ref{prop:ii.1.39}e) 以及公式 (33) 推出.
\end{theorem}

当 $r=-2$ 时, 定理中的映射成为
\[
 \Gal(L/K)^{\mathrm{ab}}\xrightarrow{\simeq}K^\times/\Nm_{L/K}L^\times.
\]
我们把它的逆映射记作
\[
 \phi_{L/K}\colon K^\times/\Nm_{L/K}L^\times\to\Gal(L/K)^{\mathrm{ab}}
\]
并称之为\emph{局部 Artin 映射}, 或者记作 $\mathrm{rec}_{L/K}$ 并称之为\emph{局部互反映射}. (后一个名字更常见, 但很难看出这个映射里有什么互反性.)

\begin{lemma}\label{lem:iii.3.2}
设 $L\supset E\supset K$ 为局部域, 且 $L/K$ Galois. 则下列图表交换:
\[
\begin{tikzcd}
	{E^\times} & {\Gal(L/E)^{\mathrm{ab}}} \\
	{K^\times} & {\Gal(L/K)^{\mathrm{ab}}}
	\arrow["{\phi_{L/E}}", from=1-1, to=1-2]
	\arrow["{\Nm_{E/K}}"', from=1-1, to=2-1]
	\arrow[from=1-2, to=2-2]
	\arrow["{\phi_{L/K}}"', from=2-1, to=2-2]
\end{tikzcd}\qquad
\begin{tikzcd}
	{E^\times} & {\Gal(L/E)^{\mathrm{ab}}} \\
	{K^\times} & {\Gal(L/K)^{\mathrm{ab}}}
	\arrow["{\phi_{L/E}}", from=1-1, to=1-2]
	\arrow[from=2-1, to=1-1]
	\arrow["{\phi_{L/K}}"', from=2-1, to=2-2]
	\arrow["{\operatorname{Ver}}", from=2-2, to=1-2]
\end{tikzcd}
\]
两图中未标名的竖直箭头分别由包含映射 $\Gal(L/E)\hookrightarrow\Gal(L/K)$ 与 $K^\times\hookrightarrow E^\times$ 诱导; 第二个图右侧的竖直箭头为转移映射 $\operatorname{Ver}$.
\end{lemma}

\begin{proof}
待补 (第二个图表其实并不需要).
\end{proof}

\noindent\textbf{3.3} 设 $L\supset E\supset K$ 为局部域, 且 $L$ 与 $E$ 都在 $K$ 上 Galois. 则下列图表交换:
\[
\begin{tikzcd}[column sep=small]
	{K^\times} & {\Gal(L/K)^{\mathrm{ab}}} \\
	& {\Gal(E/K)^{\mathrm{ab}}}
	\arrow["{\phi_{L/K}}", from=1-1, to=1-2]
	\arrow["{\phi_{E/K}}"', from=1-1, to=2-2]
	\arrow[from=1-2, to=2-2]
\end{tikzcd}
\]
未标名的箭头由满射 $\sigma\mapsto\sigma|E\colon\Gal(L/K)\to\Gal(E/K)$ 诱导.

\begin{proof}
这直接由局部 Artin 映射的定义并利用 (33) 推出.\footnote{Hervé Jacquet 指出, 这个陈述似乎并不正确, 但利用命题\ref{prop:iii.3.6} 可以证明它; 我不知道不用 \ref{prop:iii.3.6} 能否证明.}
\end{proof}

% 3.3 是不带环境名的编号条目, 不占定理计数器; 调至 3 使下一条显示为 3.4
\setcounter{theorem}{3}

特别地, 若 $L\supset E\supset K$ 是 $K$ 的有限 Abel 扩张塔, 则对所有 $a\in K^\times$ 有 $\phi_{L/K}(a)|E=\phi_{E/K}(a)$, 于是我们可以定义 $\varphi_K\colon K^\times\to\Gal(K^{\mathrm{ab}}/K)$ 为这样的同态: 对每个有限 Abel 扩张 $L/K$, $\varphi_K(a)|L=\varphi_{L/K}(a)$.

\begin{theorem}\label{thm:iii.3.4}
对每个局部域 $K$, 存在一个同态 (局部 Artin 映射)
\[
 \phi_K\colon K^\times\to\Gal(K^{\mathrm{ab}}/K),
\]
具有以下性质:
\begin{enumerate}[(a)]
\item 对 $K$ 的每个素元 $\pi$, $\phi_K(\pi)|K^{\mathrm{un}}=\Frob_K$;
\item 对 $K$ 的每个有限 Abel 扩张 $L$, $\Nm_{L/K}(L^\times)$ 包含在映射 $a\mapsto\phi_K(a)|L$ 的核中, 且 $\phi_K$ 诱导出同构
\[
 \phi_{L/K}\colon K^\times/\Nm_{L/K}(L^\times)\to\Gal(L/K).
\]
\end{enumerate}
\end{theorem}

\begin{proof}
除了 (a) 之外一切都由上述内容显然得出; 而 (a) 可由以下事实推出: 对 $K$ 的非分歧扩张 $L$, $\phi_{L/K}$ 与 \S 1 中定义的那个一致, 于是我们可应用命题\ref{prop:iii.1.9}.
\end{proof}

对读过第 \ref{ch:I} 章 \S 3 的读者来说, 这就完成了局部类域论主定理的证明 —— 如果愿意, 他们现在可以跳到第 \ref{ch:V} 章. 其他读者则需要继续读至本章结束.

\begin{theorem}[范限制定理]\label{thm:iii.3.5}
设 $L$ 为 $K$ 的有限扩张, 并设 $E$ 为 $L$ 中包含的 $K$ 的最大 Abel 扩张; 则
\[
 \Nm_{L/K}(L^\times)=\Nm_{E/K}(E^\times).
\]
\end{theorem}

\begin{proof}
由范映射的传递性, $\Nm_{L/K}(L^\times)$ 是 $\Nm_{E/K}(E^\times)$ 的子群. 若 $L/K$ 是 Galois 扩张, 则 $\Gal(E/K)=\Gal(L/K)^{\mathrm{ab}}$, 于是定理\ref{thm:iii.3.1} 表明两个范群在 $K^\times$ 中都有指数 $[E:K]$. 这意味着它们相等.

在一般情形下, 我们取 $K$ 的一个包含 $L$ 的有限 Galois 扩张 $L'$. 令 $G=\Gal(L'/K)$, $H=\Gal(L'/L)$. 于是
\[
\begin{tikzcd}
	{L'} \\
	L \\
	E \\
	K
	\arrow[from=1-1, to=2-1]
	\arrow[from=2-1, to=3-1]
	\arrow[from=3-1, to=4-1]
\end{tikzcd}
\]
那么 $E$ 是 $L'$ 中在 $K$ 上 Abel 且包含于 $L$ 的最大子域. 因此固定它的 $G$ 的子群是 $G'\cdot H$, 其中 $G'$ 是 $G$ 的派生群. 设 $a\in\Nm(E^\times)$. 我们必须证明 $a\in\Nm(L^\times)$. 考虑交换图表
\[
\begin{tikzcd}
	{L^\times} & {H/H'} \\
	{K^\times} & {G/G'} \\
	{K^\times} & {G/G'H}
	\arrow["{\phi_{L'/L}}", from=1-1, to=1-2]
	\arrow["\Nm"', from=1-1, to=2-1]
	\arrow[from=1-2, to=2-2]
	\arrow["{\phi_{L'/K}}", from=2-1, to=2-2]
	\arrow[equal, from=2-1, to=3-1]
	\arrow[from=2-2, to=3-2]
	\arrow["{\phi_{E/K}}"', from=3-1, to=3-2]
\end{tikzcd}
\]
$G/G'$ 中的元素 $\phi_{L'/K}(a)$ 在 $G/G'H$ 中映到 $1$. 由于 $\phi_{L'/L}$ 是满射, 我们从图表中看出存在 $b\in L^\times$ 使得 $\phi_{L'/K}(a)=\phi_{L'/K}(\Nm(b))$, 从而 $a/\Nm(b)\in\Nm(L^\times)$, 设 $a/\Nm(b)=\Nm(c)$. 现在
\[
 a=\Nm_{L/K}(b\cdot\Nm_{L'/L}(c))\in\Nm_{L/K}(L^\times).
\]
\end{proof}

定理\ref{thm:iii.3.5} 表明, 想用范群来分类局部域的非 Abel 扩张是没有希望的.

\subsubsection*{局部 Artin 映射的另一描述}

设 $L/K$ 为 Galois 群为 $G$ 的有限 Abel 扩张, 并设 $u_{L/K}\in H^2(G,L^\times)$ 为基本类. 局部 Artin 映射 $\phi_{L/K}$ 是同构
\[
 x\mapsto x\cup u_{L/K}\colon H_T^{-2}(G,\Z)\to H_T^0(G,L^\times)
\]
的逆. 这个定义难以操作, 因为同时涉及同调群与上同调群的杯积没有很方便的描述. 我们改为纯粹用上同调群来重新解释这个映射. 考虑杯积配对
\[
 H^0(G,L^\times)\times H^2(G,\Z)\ \longrightarrow\ H^2(G,L^\times)\xrightarrow{\inv_{L/K}}\Q/\Z.
\]
给定元素 $a\in H^0(G,L^\times)=K^\times$ 以及由上循环 $f\colon G\times G\to\Z$ 表示的类 $c\in H^2(G,\Z)$, 杯积类 $a\cup c$ 由上循环 $(\sigma,\tau)\mapsto a^{f(\sigma,\tau)}$ 表示. 还要回顾我们有一个同构
\[
 \Hom(G,\Q/\Z)=H^1(G,\Q/\Z)\xrightarrow{\delta}H^2(G,\Z).
\]

\begin{proposition}\label{prop:iii.3.6}
对每个 $\chi\in\Hom_{\mathrm{cts}}(G,\Q/\Z)$ 与 $a\in K^\times$,
\[
 \chi(\phi_{L/K}(a))=\inv_K(a\cup\delta\chi).
\]
\end{proposition}

\begin{proof}
见 \cite{Serre1962}, 第 XI 章的 ``附录'', 以及 \cite{Serre1967a}, p. 140.
\end{proof}

利用这一点, 我们可以给出命题\ref{thm:iii.3.4} 的另一个证明.

\begin{lemma}\label{lem:iii.3.7}
若 $L/K$ 非分歧, 则 $\phi_{L/K}$ 把 $a\in K^\times$ 送到 $\Frob^{\ord_K(a)}$.
\end{lemma}

\begin{proof}
回顾 $\inv_K$ 定义为复合映射
\[
 H^2(G,L^\times)\xrightarrow{\ord}H^2(G,\Z)\xleftarrow{\delta}\Hom(G,\Q/\Z)\xrightarrow{\chi\mapsto\chi(\sigma)}\Q/\Z.
\]
由杯积的函子性
\[
 \ord(a\cup\delta\chi)=\ord(a)\cup\delta\chi,\quad a\in K^\times,\quad \chi\in\Hom(G,\Q/\Z),
\]
这里左边的 $\ord$ 表示由 $\ord_L\colon L^\times\to\Z$ 在 $H^2$ 上诱导的映射, 右边则是该映射本身. 设 $a\in H^0(G,L^\times)=K^\times$, 并令 $m=\ord_L(a)$. 对每个 $\chi\in\Hom(G,\Q/\Z)$, 在上面的图表中
\[
 a\cup\delta\chi\mapsto\ord(a)\cup\delta\chi\mapsto m\chi\mapsto\chi(\sigma^m),\quad \sigma=\Frob,
\]
即 $\inv_K(a\cup\delta\chi)=\chi(\sigma^m)$. 把它与 (\ref{prop:iii.3.6}) 中的公式结合起来, 我们得到
\[
 \chi(\phi_{L/K}(a))=\chi(\sigma^{\ord_K(a)})
\]
对所有 $\chi\in\Hom(G,\Q/\Z)$ 成立, 从而 $\phi_{L/K}(a)=\sigma^{\ord_K(a)}$.
\end{proof}

对 $G$ 的每个特征标 $\chi$, 我们得到 $K^\times$ 的一个特征标 $a\mapsto\inv_K(a\cup\delta\chi)$. 由对偶性我们得到一个映射 $K^\times\to G$. 命题\ref{prop:iii.3.6} 表明这个映射就是 $\phi_{L/K}\colon K^\times\to\Gal(L/K)$.

% ==================================================================
\section{Hilbert 符号}\label{sec:III.4}\index{hilbertfuhao@Hilbert 符号(Hilbert symbol)}

由于 $K^\times$ 中有限指数子群的交是 $\{1\}$, 存在定理蕴含范群的交为 $1$. 下述关键结果是我们证明存在定理的出发点.

\begin{proposition}\label{prop:iii.4.1}
设 $K$ 为一个含有 $n$ 次本原单位根的局部域. $K^\times$ 中任何元素, 若它是 $K$ 的每个次数整除 $n$ 的循环扩张的范, 则它是一个 $n$ 次幂.
\end{proposition}

这将由 Hilbert 符号的理论推出.

\subsection*{一个非常特殊的情形}

在这一小节中, 我们说明如何只用 \cite{Serre1970} 第 III 章中的初等结果来得到 $K=\Q_p$, $n=2$ 这一特殊情形下的命题.

对 $a,b\in\Q_p^\times$, 定义
\[
 (a,b)_p=
 \begin{cases}
 \hphantom{-}1 & \text{若 } z^2=ax^2+by^2 \text{ 在 } \Q_p \text{ 中有非平凡解},\\
 -1 & \text{其他情形}.
 \end{cases}
\]
显然 $(a,b)_p$ 只依赖于 $a$ 与 $b$ 模平方的类, 于是这是一个配对
\[
 a,b\mapsto(a,b)_p\colon \Q_p^\times/\Q_p^{\times 2}\times\Q_p^\times/\Q_p^{\times 2}\to\{\pm 1\}.
\]
Serre 证明了 (同上, 1.1, 1.2) 这个配对是双乘性的
\[
 (aa',b)_p=(a,b)_p(a',b)_p,\quad (a,bb')_p=(a,b)_p(a,b')_p,
\]
对称的 (在这种情形下与斜对称是一回事)
\[
 (b,a)_p=(a,b)_p^{-1}=(a,b)_p,
\]
并且是非退化的 (即左核与右核都是平凡的).

设 $a$ 为 $\Q_p^\times$ 中一个非平方元. 那么
\begin{align*}
 &b \text{ 是来自 } \Q_p[\sqrt{a}] \text{ 的范}\\
 \iff{}& b=z^2-ax^2 \text{ 在 } \Q_p \text{ 中有解}\\
 \iff{}& z^2=ax^2+by^2 \text{ 在 } \Q_p \text{ 中有非平凡解}\\
 \iff{}& (a,b)_p=1.
\end{align*}
于是, 由配对的退化性,
\[
 b \text{ 对所有 } a \text{ 都是来自 } \Q_p[\sqrt{a}] \text{ 的范}\implies (a,b)_p=1 \text{ 对所有 } a \implies b\in\Q_p^{\times 2}.
\]
这就在此情形下证明了命题.

\begin{aside}\label{aside:iii.4.2}
映射 $a,b\mapsto(a,b)_p$ 称为 Hilbert 符号. 它对 $\R$ 也有定义, 并且 Serre 证明了 (Hilbert 乘积公式)
\begin{equation}
 \prod_{p\le\infty}(a,b)_p=1.\tag{34}
\end{equation}
\end{aside}

\subsubsection*{一般情形}

在这一小节中, $\Hom(G,\Z/n\Z)$ 表示连续同态 $G\to\Z/n\Z$ 构成的群.

\paragraph*{步骤 1: 回顾.} 回顾对每个循环群 $\bar G$ 与 $\bar G$-模 $M$, 都有同构
\[
 H_T^r(\bar G,M)\to H_T^{r+2}(\bar G,M),
 \]
一旦选定 $\bar G$ 的一个生成元 $\sigma$, 它们就成为典范的. 更确切地说, 令 $n=(\bar G:1)$, 并设 $\chi$ 为把 $\sigma$ 送到 $1\mod n$ 的同构 $\bar G\to\Z/n\Z$. 从
\[
 0\to\Z\xrightarrow{n}\Z\to\Z/n\Z\to 0
\]
的上同调序列, 我们得到一个边缘映射 $\delta\colon H^1(\bar G,\Z/n\Z)\to H^2(\bar G,\Z)$. 同构
\[
 H_T^r(\bar G,M)\to H_T^{r+2}(\bar G,M)
\]
就是 $b\mapsto\delta\chi\cup b$.

现在设 $K$ 为一个域 (不必是局部域), 并设 $L/K$ 为 Galois 群为 $\bar G$ (阶为 $n$) 的循环扩张. 当 $r=0$ 且 $M=L^\times$ 时, 该同构成为
\[
 b\mapsto\delta\chi\cup b\colon K^\times/\Nm L^\times\to H^2(\bar G,L^\times).
\]
对 $\chi\in\Hom(\bar G,\Z/n\Z)$ 与 $b\in K^\times$, 定义
\[
 (\chi,b)'=\delta\chi\cup b\in H^2(\bar G,L^\times).
\]
上面的讨论表明, 若 $\chi$ 是同构 $\bar G\to\Z/n\Z$, 则
\[
 (\chi,b)'=0\iff b \text{ 是来自 } L^\times \text{ 的范}.
\]

\paragraph*{步骤 2.} 设 $K$ 为一个域 (不必是局部域). 若 $K$ 的特征为 $p\neq 0$, 则我们要求 $n$ 与 $p$ 互素. 设 $K^{\mathrm{al}}$ 为 $K$ 的一个可分代数闭包 (于是当 $K$ 完美时它就是一个代数闭包), 并设 $G=\Gal(K^{\mathrm{al}}/K)$. 设 $\mu_n$ 为 $K^{\mathrm{al}}$ 中 $n$ 次单位根构成的群, 视为 $G$-模. 在我们对 $n$ 的假设下, $\mu_n$ 是一个 $n$ 阶循环群.

从
\[
 0\to\mu_n\to K^{\mathrm{al}\times}\xrightarrow{x\mapsto x^n}K^{\mathrm{al}\times}\to 0
\]
的上同调序列以及 Hilbert 定理 90 (II, \ref{prop:ii.1.22}), 我们得到
\[
 \delta\colon K^\times/K^{\times n}\xrightarrow{\simeq}H^1(G,\mu_n),\qquad H^2(G,\mu_n)\xrightarrow{\simeq}H^2(G,K^{\mathrm{al}\times})_n.
\]
这里 $H^2(G,K^{\mathrm{al}\times})_n$ 表示 $H^2(G,K^{\mathrm{al}\times})$ 中被 $n$ 杀死的元素构成的子群.

考虑杯积配对
\[
\begin{tikzcd}[column sep=small]
	{H^1(G,\Z/n\Z)} & {\times} & {H^1(G,\mu_n)} & {\longrightarrow} & {H^2(G,\mu_n)} \\
	{\Hom(G,\Z/n\Z)} & & {K^\times/K^{\times n}} & & {H^2(G,K^{\mathrm{al}\times})_n}
	\arrow[equal, no head, from=1-1, to=2-1]
	\arrow[equal, no head, from=1-3, to=2-3]
	\arrow[equal, no head, from=1-5, to=2-5]
\end{tikzcd}
\]
对 $\chi\in\Hom(G,\Z/n\Z)$ 与 $b\in K^\times/K^{\times n}$, 我们把 $(\chi,b)$ 记作该偶对在上述配对下的像: $(\chi,b)=\chi\cup\delta b$ (作为 $H^2(G,K^{\mathrm{al}\times})$ 中的元素).

设 $\chi\in\Hom(G,\Z/n\Z)$, 并 (为简单起见) 假设 $\chi$ 的阶为 $n$. 设 $L_\chi$ 为 $K^{\mathrm{al}}$ 中被 $\Ker(\chi)$ 固定的子域. 于是 $L_\chi$ 是 $K$ 的 $n$ 次循环扩张, 且 $\chi$ 诱导出其 Galois 群 $\bar G$ 与 $\Z/n\Z$ 的一个同构. 设 $b\in K^\times$, 并令
\[
 (\chi,b)'=\delta\chi\cup b\in H^2(\bar G,L^\times),
\]
如同步骤 1. 那么
\[
 \Inf((\chi,b)')=(\chi,b).
\]
由于膨胀映射 $H^2(\bar G,L_\chi)\to H^2(G,K^{\mathrm{al}\times})$ 是单射 (见 II, \ref{ex:ii.1.36}), 我们推出
\[
 (\chi,b)=0\iff b \text{ 是来自 } L_\chi^\times \text{ 的范}.
\]

\paragraph*{步骤 3.} 现在假设 $K$ 是局部域, 于是我们有一个典范同构
\[
 \inv_K\colon H^2(G,K^{\mathrm{al}\times})\to\Q/\Z.
\]
对 $\chi\in\Hom(G,\Z/n\Z)\simeq H^1(G,\Z/n\Z)$ 与 $b\in K^\times/K^{\times n}\simeq H^1(G,\mu_n)$, 我们现在定义
\[
 (\chi,b)=\inv_K(\chi\cup\delta b)\in\tfrac{1}{n}\Z/\Z.
\]
于是我们有一个典范配对
\[
 \chi,b\mapsto(\chi,b)\colon H^1(G,\Z/n\Z)\times H^1(G,\mu_n)\to H^2(G,\mu_n)\simeq\tfrac{1}{n}\Z/\Z.
\]

\begin{proposition}\label{prop:iii.4.3}
这个配对的左核为零.
\end{proposition}

\begin{proof}
设 $\chi\in\Hom(G,\Z/n\Z)$; 我们必须证明: 若对所有 $b\in K^\times/K^{\times n}$ 都有 $(\chi,b)=0$, 则 $\chi=0$.

回顾局部 Artin 映射 $\phi_K\colon K^\times\to\Gal(K^{\mathrm{ab}}/K)$ 与
\[
 \tau\mapsto\tau|L\colon \Gal(K^{\mathrm{ab}}/K)\to\Gal(L/K)
\]
的复合对每个包含在 $K^{\mathrm{ab}}$ 中的 $K$ 的有限扩张 $L$ 都是满射. 这蕴含 $\phi_K$ 的像在 $\Gal(K^{\mathrm{ab}}/K)$ 中稠密. 回顾 (\ref{prop:iii.3.6}): 对 $b\in K^\times$,
\[
 (\chi,b\mod K^{\times n})=\chi(\phi_K(b)).
\]
因为 $\chi$ 连续, 若它在 $\phi_K(K^\times)$ 上为零, 则它在整个 $\Gal(K^{\mathrm{ab}}/K)$ 上为零.
\end{proof}

\paragraph*{步骤 4.} 现在假设 $K$ 含有一个 $n$ 次本原单位根. 于是作为 $G$-模有 $\mu_n\approx\Z/n\Z$ (该同构依赖于一个 $n$ 次本原单位根的选取). 特别地, $\mu_n\otimes\mu_n\approx\mu_n$, 从而 $H^2(G,\mu_n\otimes\mu_n)\approx H^2(G,\mu_n)$ (非典范). 存在一个典范同构
\[
 H^2(G,\mu_n)\otimes\mu_n\xrightarrow{\simeq}H^2(G,\mu_n\otimes\mu_n),
\]
它可以定义为 $x,y\mapsto x\cup y$. 把它与由不变映射给出的同构 $H^2(G,\mu_n)\to\tfrac{1}{n}\Z/\Z$ 复合, 我们得到一个同构
\[
 H^2(G,\mu_n\otimes\mu_n)\simeq\mu_n.
\]
于是我们有一个杯积配对
\[
\begin{tikzcd}[column sep=small]
	{H^1(G,\mu_n)} & {\times} & {H^1(G,\mu_n)} & {\longrightarrow} & {H^2(G,\mu_n\otimes\mu_n)} \\
	{K^\times/K^{\times n}} & & {K^\times/K^{\times n}} & & {\mu_n}
	\arrow[equal, no head, from=1-1, to=2-1]
	\arrow[equal, no head, from=1-3, to=2-3]
	\arrow[equal, no head, from=1-5, to=2-5]
\end{tikzcd}
\]
对 $a,b\in K^\times$, 它们在这个配对下于 $\mu_n$ 中的像记作 $(a,b)$. 这个配对称为 \emph{Hilbert 符号}.

\begin{theorem}\label{thm:iii.4.4}
Hilbert 符号具有下列性质.
\begin{enumerate}[(a)]
\item 它是双乘性的, 即
\[
 (aa',b)=(a,b)(a',b),\quad (a,bb')=(a,b)(a,b').
\]
\item 它是斜对称的, 即
\[
 (b,a)=(a,b)^{-1}.
\]
\item 它是非退化的, 即
\[
 (a,b)=1 \text{ 对所有 } b\in K^\times/K^{\times n}\implies a\in K^{\times n},
\]
\[
 (a,b)=1 \text{ 对所有 } a\in K^\times/K^{\times n}\implies b\in K^{\times n}.
\]
\item $(a,b)=1$ 当且仅当 $b$ 是来自 $K[\sqrt[n]{a}]$ 的范.
\end{enumerate}
\end{theorem}

\begin{proof}
陈述 (a) 与 (b) 可由 Hilbert 符号用杯积的定义直接得出.

接下来证明 (c). 选取一个 $n$ 次单位根, 并用它来定义同构 $\Z/n\Z\to\mu_n$ 与 $\mu_n\to\mu_n\otimes\mu_n$. 由这些映射定义的图表
\[
\begin{tikzcd}[column sep=small, row sep=large]
	{H^1(G,\mu_n)} & {\times} & {H^1(G,\mu_n)} & {\longrightarrow} & {H^2(G,\mu_n\otimes\mu_n)} & {\simeq} & {H^2(G,\mu_n)\otimes\mu_n} \\
	{H^1(G,\Z/n\Z)} & {\times} & {H^1(G,\mu_n)} & {\longrightarrow} & {H^2(G,\mu_n)}
	\arrow[from=2-1, to=1-1]
	\arrow[equal, no head, from=2-3, to=1-3]
	\arrow[from=2-5, to=1-5]
\end{tikzcd}
\]
交换. 由命题\ref{prop:iii.4.3}, 下方配对的左核为零. 因此上方配对的左核也为零, 而由斜对称性, 这蕴含它的右核为零.

最后证明 (d). 假设 (为简单起见) $K[\sqrt[n]{a}]$ 在 $K$ 上的次数为 $n$. 取一个 $n$ 次本原单位根 $\zeta$. 那么 $K[\sqrt[n]{a}]$ 在 $K$ 上 $n$ 次循环, 其 Galois 群由把 $\sqrt[n]{a}$ 送到 $\zeta\sqrt[n]{a}$ 的映射 $\sigma$ 生成. 取定这个 $\zeta$ 之后, $a\in K^\times/K^{\times n}$ 对应于把 $\sigma$ 送到 $1\mod n\Z$ 的元素 $\chi_a\in\Hom(G,\Z/n\Z)$. 由于 $K[\sqrt[n]{a}]=L_{\chi_a}$, (d) 由上面的公式推出.
\end{proof}

\begin{proof}[命题\ref{prop:iii.4.1} 的证明]
我们可以像 $K=\Q_p$, $n=2$ 的情形那样, 从定理\ref{thm:iii.4.4} 推出命题\ref{prop:iii.4.1}: 若 $b$ 对所有 $a$ 都是来自 $K[\sqrt[n]{a}]$ 的范, 则 (由 \ref{thm:iii.4.4}d) 对所有 $a$ 有 $(a,b)=1$, 从而 (由 \ref{thm:iii.4.4}c) $b\in K^{\times n}$.
\end{proof}

\begin{remark}\label{rem:iii.4.5}
Hilbert 符号通过公式
\[
 \phi_K(b)(a^{\frac{1}{n}})=(a,b)a^{\frac{1}{n}}
\]
与局部 Artin 映射相联系. 注意 Galois 理论告诉我们, 对每个 $\tau\in\Gal(K[a^{\frac{1}{n}}]/K)$, 都有 $\tau a^{\frac{1}{n}}=\zeta a^{\frac{1}{n}}$, 其中 $\zeta$ 是某个 $n$ 次单位根 (记住, 我们假设 $K$ 含有 $n$ 次单位根), 所以这个公式的要点在于单位根是相同的. 该公式的证明是一个从命题\ref{prop:iii.3.6} 出发的杯积练习.
\end{remark}

\begin{remark}\label{rem:iii.4.6}
存在定理将表明 $\phi_K$ 定义了一个同构 $K^\times/K^{\times n}\to\Gal(L/K)$, 其中 $L/K$ 是 $K$ 的指数为 $n$ 的最大 Abel 扩张 (没有存在定理时我们只知道这个映射是满射). 显然 $\Hom(G,\Z/n\Z)=\Hom(\Gal(L/K),\Z/n\Z)$, 后者是 $\Gal(L/K)$ 的对偶. 于是 $K^\times/K^{\times n}$ 与 $\Hom(G,\Z/n\Z)$ 有相同的阶. 这表明步骤 2 中的配对
\[
 H^1(G,\Z/n\Z)\times H^1(G,\mu_n)\to H^2(G,\mu_n)\simeq\tfrac{1}{n}\Z/\Z
\]
是非退化的. 更一般地, 对每个有限 $G$-模 $M$, 杯积配对
\[
 H^r(G,M)\times H^{2-r}(G,\check M)\to H^2(G,\mu_n)\simeq\tfrac{1}{n}\Z/\Z
\]
都是非退化的 —— 这里 $\check M=\Hom(M,\mu_n)$. 这个对偶定理由 Poitou 与 Tate 在 1950 年代早期证明, 它是我的书 ``Arithmetic Duality Theorems'' 中那些定理的一个出发点. 注意它蕴含 $H^2(G,\mu_n\otimes\mu_n)\simeq(\mu_n)_G$.
\end{remark}

\begin{remark}\label{rem:iii.4.7}
Hilbert 符号可以用中心单代数来解释 (见下一章). 对 $a,b\in K^\times$ 以及 $n$ 次本原单位根 $\zeta$, 定义 $A(a,b;\zeta)$ 为以 $i$, $j$ 为生成元、满足关系 $i^n=a$, $j^n=b$, $ij=\zeta ji$ 的 $K$-代数. 那么 $(a,b)$ 就是 $A(a,b;\zeta)$ 所代表的 $\mathrm{Br}(K)_n$ 中的类. 注意当 $n=2$ 时我们重新得到四元数代数 $H(a,b)$.
\end{remark}

\begin{remark}\label{rem:iii.4.8}
Hilbert 符号
\[
 (\ ,\ )\colon K^\times/K^{\times n}\times K^\times/K^{\times n}\to\mu_n
\]
对 $K=\R$ 与 $n=2$ 也有定义. 此外, 乘积公式 (34) 对所有数域与所有 $n$ 都成立 (见下文 V, \ref{ex:v.5.4}).
\end{remark}

% ==================================================================
\section{存在定理}\label{sec:III.5}\index{cunzaidingli@存在定理(existence theorem)}

设 $K$ 为局部域. 回顾 $K^\times$ 的子群 $N$ 称为一个\emph{范群}, 如果存在有限 Abel 扩张 $L/K$ 使得 $\Nm_{L/K}(L^\times)=N$. 由于 $K^\times/N\xrightarrow{\simeq}\Gal(L/K)$, 这样的群 $N$ 在 $K^\times$ 中指数有限, 从而是开的 (I, \ref{lem:1.3}).

\begin{theorem}[存在定理]\label{thm:iii.5.1}
$K^\times$ 的每个有限指数的开子群都是范群.
\end{theorem}

在这份草稿中, 我只对特征零的 $K$ 证明这一点. 此时有限指数的子群是开的, 于是我们可以使用 Hilbert 符号的理论.

\begin{remark}\label{rem:iii.5.2}
在开始定理\ref{thm:iii.5.1} 的证明之前, 我们应当考虑一下抽象的情形: 设 $Z$ 为一个无限 Abel 群, 并设 $\mathcal{N}$ 为一族有限指数的子群; 为证明 $\mathcal{N}$ 包含 $Z$ 的所有有限指数子群, 我们需要知道些什么? 显然, 一个起点是要知道:

(a) 若 $I$ 包含 $\mathcal{N}$ 的一个元素, 则 $I$ 在 $\mathcal{N}$ 中;

(b) 若 $N_1$ 与 $N_2$ 在 $\mathcal{N}$ 中, 则 $N_1\cap N_2$ 包含 $\mathcal{N}$ 的一个元素 (从而在 $\mathcal{N}$ 中). 我们知道 $K^\times$ 的范子群族具有这些性质 (推论\ref{prop:iii.1.2}), 但是例如 $\Z$ 中包含 $n\Z$ 的子群构成的集合也具有这些性质. 显然我们需要知道更多. 设 $D$ 为 $\mathcal{N}$ 中所有群的交. 若 $Z$ 的每个有限指数子群 $I$ 都包含 $\mathcal{N}$ 的一个元素, 则它包含 $D$. 反过来, 若 $D$ 是可除的, 即对所有整数 $n$ 有 $D=nD$, 那么 $D$ 包含在每个有限指数的子群中 (指数为 $n$ 的 $Z$ 的子群包含 $nZ$, 而 $nZ$ 包含 $nD=D$). 因此我们需要 $D$ 可除, 但即使这样也还不够. 例如, 取 $Z=\Z$ 与 $\mathcal{N}=\{p^n\Z\mid n\in\N\}$. 那么 $D=\bigcap p^n\Z=0$ 是可除的, 但并非 $\Z$ 的每个有限指数子群都包含 $\mathcal{N}$ 的某个元素. [也许我应当列出一族 $\mathcal{N}$ 为包含所有有限指数的开子群所应满足的完整 ``公理'' 清单 —— 见 \cite{Serre1962}, XI.5, 或 Artin 与 Tate \citeyear{ArtinTate1961}, 第 14 章. 这对我们处理整体存在定理时或许有用.]
\end{remark}

\subsubsection*{定理\ref{thm:iii.5.1} 的证明}

\paragraph*{步骤 1.} 对所有有限扩张 $L/K$, 范映射 $L^\times\to K^\times$ 有闭的像与紧的核.

\begin{proof}
由于像的指数有限, 它是开的 (见 I, \ref{lem:1.3}), 从而是闭的. 核是闭的, 且等式
\[
 \ord_K(\Nm(a))=f\cdot\ord_L(a)
\]
表明它包含在 $U_L$ 中, 而 $U_L$ 是紧的.
\end{proof}

令
\[
 D_K=\bigcap \Nm_{L/K}(L^\times),
\]
其中 $L$ 跑遍 $K$ 的有限扩张 (或者只跑遍有限 Abel 扩张 —— 由定理\ref{thm:iii.3.5}, 这是同一个群).

\paragraph*{步骤 2.} 对任意有限扩张 $K'/K$, 有 $\Nm_{K'/K}D_{K'}=D_K$.

\begin{proof}
设 $a\in D_K$, 并考虑诸集合
\[
 \Nm_{L/K'}L^\times\cap\Nm_{K'/K}^{-1}(a)
\]
其中 $L/K'$ 有限. 它们是紧的且非空的, 并且任意两个都包含第三个. 因此它们的交非空 (点集拓扑). 交中的一个元素落在 $D_{K'}$ 中且范为 $a$.
\end{proof}

\paragraph*{步骤 3.} 群 $D_K$ 是可除的.

\begin{proof}
设 $n>1$ 为一个整数 —— 我们必须证明 $nD_K=D_K$. 设 $a\in D_K$. 对每个含有 $n$ 次本原单位根的 $K$ 的有限扩张 $L$, 考虑集合
\[
 E(L)=\{b\in K^\times\mid b^n=a,\ b\in\Nm_{L/K}L^\times\}.
\]
它是非空的: 对某个 $a'\in D_L$ 有 $a=\Nm_{L/K}a'$; 依命题\ref{prop:iii.4.1}, 对某个 $c\in L$ 有 $a'=c^n$, 且
\[
 \Nm_{L/K}(c)^n=\Nm_{L/K}(a')=a;
\]
因此 $b\stackrel{\mathrm{def}}{=}\Nm_{L/K}(c)\in E(L)$. 此外 $E(L\cdot L')\subset E(L)\cap E(L')$. 由于每个集合都是有限的, 它们的交非空. 交中的一个元素落在 $D_K$ 中且 $n$ 次幂为 $a$.
\end{proof}

\paragraph*{步骤 4.} $D_K=\{1\}$. [实际上, 我们这里不用这一点.]

\begin{proof}
取 $K$ 的一个素元. 令 $V_{m,n}=U^{(m)}\times\pi^{n\Z}$. 那么 $V_{m,n}$ 是 $K^\times$ 的一个 (开的) 有限指数子群, 因此包含 $D_K$. 现在 $D_K\subset\bigcap_{m,n}V_{m,n}=\{1\}$.
\end{proof}

\paragraph*{步骤 5.} $K^\times$ 中每个包含 $U_K$ 的有限指数子群 $I$ 都是范子群.

\begin{proof}
由于 $\ord_K\colon K^\times\to\Z$ 是满射且核为 $U_K$, 所说的那些子群 $I$ 就是 $K^\times$ 中形如 $\ord_K^{-1}(n\Z)$ ($n\ge 1$) 的子群. 设 $K_n$ 为 $K$ 的 $n$ 次非分歧扩张. 那么 $\Nm_{K_n/K}(K_n^\times)$ 是 $K^\times$ 中包含 $U_K$ (见 \ref{prop:iii.1.2}) 且在 $\Z$ 中的像为 $n\Z$ 的子群. 结论成立.
\end{proof}

我们现在来证明定理. 设 $\mathcal{N}$ 为 $K^\times$ 中范群构成的集合, 于是 $D_K=\bigcap_{N\in\mathcal{N}}N$. 设 $I$ 为 $K^\times$ 的一个有限指数子群. 因为 $D_K$ 可除, 有 $I\supset D_K$, 于是 $I\supset\bigcap_{N\in\mathcal{N}}N\supset\bigcap_{N\in\mathcal{N}}(N\cap U_K)$. 因此诸集合 $(N\cap U_K)\setminus I$ 的交为空. 由于它们是紧的, 某个有限子族的交已经为空, 即存在 $\mathcal{N}$ 的某个有限子集 $\mathcal{S}$ 使得 $I\supset\bigcap_{N\in\mathcal{S}}(N\cap U_K)$. 由于诸集合 $N\cap U_K$ 中任意两个都包含第三个, 这蕴含对某个 $N$ 有 $I\supset N\cap U_K$.

取定一个满足 $N\cap U_K\subset I$ 的范群 $N$. 我断言 $I$ 包含
\[
 N\cap\bigl(U_K\cdot(N\cap I)\bigr).
\]
事实上, 交中的每个元素都有形式 $ab$, 其中 $a\in U_K$, $b\in N\cap I$, $ab\in N$. 最后两条陈述蕴含 $a\in N$. 因此 $a\in N\cap U_K\subset I$, 从而 $ab\in I$.

现在 $N\cap I$ 在 $K^\times$ 中指数有限, 因为 $N$ 与 $I$ 都是如此, 且 $K^\times/(N\cap I)$ 单射入 $(K^\times/N)\times(K^\times/I)$. 因此 $U_K\cdot(N\cap I)$ 是 $K^\times$ 中包含 $U_K$ 的有限指数子群, 从而是一个范群 (步骤 5). 现在 $N\cap\bigl(U_K\cdot(N\cap I)\bigr)$ 作为两个范群的交, 包含一个范群. 因此 $I$ 包含一个范群, 而这蕴含它本身是一个范群.

\begin{NOTES}
由 Krasner 引理 (ANT, 7.60) 推出, 局部域的每个有限 Abel 扩张都来自整体域的某个有限 Abel 扩张的完备化. 在 1930 年代, Hasse 与 F. K. Schmidt\footnote{H. Hasse, \emph{Die Normenresttheorie relative-Abelscher Zahlkörper als Klassenkörper im Kleinen}, J. für Mathematik (Crelle) 162 (1930), 145–154. F. K. Schmidt, \emph{Zur Klassenkörpertheorie im Kleinen}, 同上, 155–168.} 能够从整体类域论的主定理推出局部类域论的主定理.

从现代的视角看, 这似乎是一种奇怪的做法. 在 1940 年代, Chevalley 在他的类域论的代数进路中直接发展了局部类域论, 并用它来构造整体类域论. F. K. Schmidt 也证明了局部类域论可以独立于整体类域论来建立.

当时, 局部 Artin 映射还没有好的描述, 也没有构造局部域的最大 Abel 扩张的明确方法 (当然 $\Q_p$ 除外).

1958 年 Dwork 给出了局部 Artin 映射的一个明确描述, 它复现在 \cite{Serre1962} 中, 但不太令人满意.

1965 年 Lubin 与 Tate 引入了 Lubin-Tate 形式群律, 给出了 $K^{\mathrm{ab}}$ 的一个明确构造以及局部 Artin 映射的一个明确描述. 不过, 他们在证明中用到了局部 Artin 映射的存在性 (我们的定理 I, \ref{thm:1.1}).

1980 年代早期, Lubin, Gold 与 Rosen 各自独立地给出了 $K^{\mathrm{ab}}=K_\pi\cdot K^{\mathrm{un}}$ 的 ``初等'' 证明. Iwasawa 在他的书 (\cite{Iwasawa1986}) 中从 Lubin-Tate 的视角展开了整个局部类域论, 也给出了 Hilbert 符号等的明确公式 (由 de Shalit 与 Wiles 得到).

其他的非上同调进路可以在 \cite{Hazewinkel1975}, \cite{Neukirch1986}, 以及 Fesenko 与 Vostokov \citeyear{FesenkoVostokov1993} 中找到. 非上同调进路的缺点自然在于它们不提供关于局部域上同调群的信息, 而这些群在其他课题 (例如椭圆曲线) 中有重要应用.

关于顺分歧扩张情形中局部 Artin 映射的明确描述, 见 Newton, Rachel, Explicit local reciprocity for tame extensions. Math. Proc. Cambridge Philos. Soc. 152 (2012), no. 3, 425--454.

在本章中, 我在很大程度上遵循了 \cite{Serre1962} 与 \cite{Serre1967a}.
\end{NOTES}
